% Options for packages loaded elsewhere
\PassOptionsToPackage{unicode}{hyperref}
\PassOptionsToPackage{hyphens}{url}
\documentclass[
  12pt,
]{ctexart}
\usepackage{xcolor}
\usepackage{amsmath,amssymb}
\setcounter{secnumdepth}{-\maxdimen} % remove section numbering
\usepackage{iftex}
\ifPDFTeX
  \usepackage[T1]{fontenc}
  \usepackage[utf8]{inputenc}
  \usepackage{textcomp} % provide euro and other symbols
\else % if luatex or xetex
  \usepackage{unicode-math} % this also loads fontspec
  \defaultfontfeatures{Scale=MatchLowercase}
  \defaultfontfeatures[\rmfamily]{Ligatures=TeX,Scale=1}
\fi
\usepackage{lmodern}
\ifPDFTeX\else
  % xetex/luatex font selection
\fi
% Use upquote if available, for straight quotes in verbatim environments
\IfFileExists{upquote.sty}{\usepackage{upquote}}{}
\IfFileExists{microtype.sty}{% use microtype if available
  \usepackage[]{microtype}
  \UseMicrotypeSet[protrusion]{basicmath} % disable protrusion for tt fonts
}{}
\makeatletter
\@ifundefined{KOMAClassName}{% if non-KOMA class
  \IfFileExists{parskip.sty}{%
    \usepackage{parskip}
  }{% else
    \setlength{\parindent}{0pt}
    \setlength{\parskip}{6pt plus 2pt minus 1pt}}
}{% if KOMA class
  \KOMAoptions{parskip=half}}
\makeatother
\usepackage{longtable,booktabs,array}
\usepackage{calc} % for calculating minipage widths
% Correct order of tables after \paragraph or \subparagraph
\usepackage{etoolbox}
\makeatletter
\patchcmd\longtable{\par}{\if@noskipsec\mbox{}\fi\par}{}{}
\makeatother
% Allow footnotes in longtable head/foot
\IfFileExists{footnotehyper.sty}{\usepackage{footnotehyper}}{\usepackage{footnote}}
\makesavenoteenv{longtable}
\setlength{\emergencystretch}{3em} % prevent overfull lines
\providecommand{\tightlist}{%
  \setlength{\itemsep}{0pt}\setlength{\parskip}{0pt}}
\usepackage{bookmark}
\IfFileExists{xurl.sty}{\usepackage{xurl}}{} % add URL line breaks if available
\urlstyle{same}
\hypersetup{
  hidelinks,
  pdfcreator={LaTeX via pandoc}}

\author{SCX}
\date{2026}

\begin{document}

\section{Rigorous Mathematical Analysis of A2}\label{rigorous-mathematical-analysis-of-a2}

\begin{quote}
Trigger: hostile review points out that \(\exp(-2M\Delta^2/(1+(M-1)\rho))\) is not a rigorous theorem. Conclusion: correct. Plugging variance inflation into Hoeffding is heuristic, not a rigorous bound. But A2 itself has a stronger argument.
\end{quote}

\begin{center}\rule{0.5\linewidth}{0.5pt}\end{center}

\subsection{\texorpdfstring{1. Is \(\exp(-2M t^2/(1+(M-1)\rho))\) a Rigorous Theorem?}{1. Is exp(-2M t\^{}2/(1+(M-1)rho)) a Rigorous Theorem?}}\label{is-exp-2m-t2-1m-1rho-a-rigorous-theorem}

\textbf{No.}

The proof of Hoeffding's inequality relies on:
\[\mathbb{E}\left[\exp\left(\lambda\sum_{i=1}^M (X_i - \mu_i)\right)\right] = \prod_{i=1}^M \mathbb{E}[\exp(\lambda(X_i - \mu_i))]\]

This \textbf{factorization} requires \(\{X_i\}\) to be independent. If \(X_i\) are correlated, the MGF cannot be factorized, and you cannot simply replace \(M\) with \(M/(1+(M-1)\rho)\) to obtain a rigorous exponential bound.

The variance inflation \((1+(M-1)\rho)\) is correct for \textbf{\(\text{Var}(\sum X_i)\)}, but not for \textbf{exponential concentration}. Variance only gives the Chebyshev bound (polynomial rate \(1/t^2\)), not an exponential bound.

\textbf{Conclusion}: Variance-inflated Hoeffding is heuristic --- widely used in survey sampling and cluster-randomized trials, but not a mathematical theorem. Including it in the SI will be caught by reviewers.

\begin{center}\rule{0.5\linewidth}{0.5pt}\end{center}

\subsection{2. Does A2 Itself Hold Up?}\label{does-a2-itself-hold-up}

\subsubsection{2.1 A Stronger Argument: A1 Implies A2}\label{a1-implies-a2}

A1 states: \(M\) experts are trained on \textbf{disjoint i.i.d. subsets}.
\[D_1, \dots, D_M \sim_{\text{i.i.d.}} P^n \quad \text{and} \quad D_i \cap D_j = \varnothing\]

Each \(f_m = \mathcal{A}(D_m)\) is a \textbf{deterministic function} of the training algorithm \(\mathcal{A}\) applied to data subset \(D_m\).

Because \(D_1, \dots, D_M\) are independent random variables, \(f_1, \dots, f_M\) are also independent random functions.

For any fixed test point \(x\):
\[e_m(x, y) = \mathbf{1}\{\ell(f_m(x), y) > \tau\}\]

Because \(f_m\) depends only on \(D_m\) and
\(D_1 \perp D_2 \perp \cdots \perp D_M\):
\[e_1 \perp e_2 \perp \cdots \perp e_M \quad \mid \quad x\]

\textbf{A2 is not an assumption --- it is a corollary of A1.}

\subsubsection{2.2 What Is the Probability Space?}\label{what-is-the-probability-space}

The probability space for Hoeffding's bound is \textbf{the randomness of the training data}, not the randomness of the test data.

\[\mathbb{P}_{D_1,\dots,D_M}\left(C(x) > \theta \mid \text{clean}, x\right) \leq \exp(-2M(\theta-\mu_s)^2)\]

This means: if you randomly partition the training data \(M\) times and independently train \(M\) experts, the resulting expert set has, with high probability, good noise-detection properties.

\textbf{After training is complete}, you have a fixed set of experts. For a given test sample \(x\), \(C(x)\) is a deterministic number. Hoeffding does not apply.

But this does not weaken the theorem --- all generalization bounds have this property. The theorem tells you ``this \textbf{method} is good,'' not ``this \textbf{specific set of experts} is good.''

\subsubsection{2.3 Why the Hostile Reviewer's Objection Is Wrong}\label{why-the-hostile-reviewers-objection-is-wrong}

\begin{quote}
``All CNNs fail together on blurry images.''
\end{quote}

This is saying: \textbf{given already-trained CNNs}, their deterministic behavior on blurry images is consistent.

But Hoeffding asks: \textbf{before training}, after randomly partitioning the data and training independently, how confident are you that you will obtain a set of experts capable of detecting noise?

The two are not in the same probability space. The hostile reviewer conflates \textbf{post-training deterministic behavior} with \textbf{pre-training probabilistic guarantee}.

\subsubsection{2.4 But a Residual Problem Remains}\label{residual-problem}

Even if A2 holds rigorously (as a corollary of A1), \textbf{estimating} \(\mu_s\) still requires clean labels. Under finite samples, the error in \(\hat{\mu}_s\) propagates into the error in \(\hat{\Delta}_s\). This is a \textbf{finite-sample estimation problem}, not an \textbf{independence assumption problem}. Corollary 4 (finite-sample correction) already handles this with conservative upper bounds.

\begin{center}\rule{0.5\linewidth}{0.5pt}\end{center}

\subsection{3. Honest Revision Suggestions}\label{honest-revision-suggestions}

\subsubsection{3.1 Do Not Introduce A2'}\label{do-not-introduce-a2-prime}

Variance-inflated Hoeffding is heuristic and not fit for inclusion in a theorem statement. Keep the original A2, but \textbf{discuss honestly in the main text and SI}.

\subsubsection{3.2 Add a Section ``On the Role of Assumption A2'' in the SI}\label{add-a-section-on-the-role-of-assumption-a2}

Contents:
1. A2 is in fact a corollary of A1 (disjoint training sets → independent experts → independent errors)
2. The probability space is training randomness, not test behavior
3. The estimation error of \(\mu_s\) under finite samples is handled independently (Corollary 4)
4. In practice, expert errors may exhibit apparent correlation (same blurry image), but this reflects feature similarity leading to similar \(\mu_s\), not statistical correlation
5. Empirical diagnosis of A2: estimate the conditional correlation coefficient of expert errors on a validation set; if it deviates significantly from zero, check whether training sets are truly disjoint

\subsubsection{3.3 Honest Wording in the Main Text}\label{honest-wording-in-main-text}

Change ``conditional independence'' to: \textgreater{} ``conditional independence of expert errors, which follows from training on disjoint data subsets (Assumption A1)''

This downgrades A2 from a ``hypothesis'' to a ``logical consequence of A1'' while noting that it depends on correct execution of A1.

\begin{center}\rule{0.5\linewidth}{0.5pt}\end{center}

\subsection{4. Final Judgment}\label{final-judgment}

\begin{longtable}[]{@{}
  >{\raggedright\arraybackslash}p{(\linewidth - 2\tabcolsep) * \real{0.6000}}
  >{\centering\arraybackslash}p{(\linewidth - 2\tabcolsep) * \real{0.4000}}@{}}
\toprule\noalign{}
\begin{minipage}[b]{\linewidth}\raggedright
Issue
\end{minipage} & \begin{minipage}[b]{\linewidth}\centering
Status
\end{minipage} \\
\midrule\noalign{}
\endhead
\bottomrule\noalign{}
\endlastfoot
Does A2 hold? & \textbf{Yes} --- A1 implies A2 (disjoint training sets → independent experts) \\
Is it violated in practice? & \textbf{Depends on A1 execution} --- if training sets are truly disjoint, A2 holds \\
Is \(\exp(-2M\Delta^2/(1+(M-1)\rho))\) rigorous? & \textbf{No} --- this is heuristic and should not enter a theorem \\
How should it be revised? & Clarify that A2 is a corollary of A1; add honest discussion in SI; do not introduce a fake A2' \\
Impact on the paper & \textbf{Low} --- no need to change theorem statements, only add a paragraph of honest discussion \\
\end{longtable}

\begin{center}\rule{0.5\linewidth}{0.5pt}\end{center}

\subsection{5. Responses to Other Hostile Review Criticisms}\label{responses-to-other-hostile-review-criticisms}

\subsubsection{``CIFAR F1=0.617, theoretical lower bound F1\(\ge\)0.976''}\label{cifar-f1-0617-theoretical-lower-bound-f1-ge-0976}

This has already been honestly handled in §4.3 (Tightness Discussion): \(\mu_s\) in the CIFAR experiment was approximately 0.45 (after 3-epoch CPU training), not the assumed 0.20. Substituting the correct \(\mu_s\) yields a lower bound of F1 \(\ge\) 0.18, consistent with the empirical value of 0.617.

\subsubsection{\texorpdfstring{``\(\mu_s\) requires clean labels → chicken-and-egg''}{``mu\_s requires clean labels → chicken-and-egg''}}\label{mu_s-requires-clean-labels-chicken-and-egg}

The cold-start protocol requires a small number of anchor samples (Corollary 1); this is a theoretically necessary cost (Thm 3 proves unidentifiability without anchors). This is not a circular definition --- it is an engineering response to the unidentifiability theorem.

\subsubsection{``Bootstrap diagnostic \(\tau\)=0.7 is arbitrary''}\label{bootstrap-diagnostic-tau-07-is-arbitrary}

Acknowledged. \(\tau=0.7\) is an initial calibration value and should be calibrated against known noisy labels in specific domains. The SI should explicitly label it as ``recommended initial value, requires domain calibration.''

\begin{center}\rule{0.5\linewidth}{0.5pt}\end{center}

\emph{Related: {[}{[}01\_noise\_detection\_guarantee{]}{]} ·
{[}{[}03\_unidentifiability{]}{]} ·
{[}{[}a2\_correlation\_analysis{]}{]}}

\end{document}
